%% =====================================================================
%%  T-09-02  Bearing Sets the Price
%%  -------------------------------------------------------------------
%%  One idea per file. Core is mandatory. Slots in canonical order.
%%  Max 700 words. No layout commands. No filler. New source material
%%  for this entry goes in sources/<BOOK>/T-09-02.tex -- never by
%%  rewriting this file (docs/RULES.md R0.1).
%%  =====================================================================
%% @kind: topic
%% @id: T-09-02
%% @title: Bearing Sets the Price
%% @subtitle: People treat you at the level you act --- until you pass the ego of whoever is above you
%% @chapter: c09
%% @order: 20
%% @status: stable
%% @sources: B3
%% @updated: 2026-09-19

\topic{T-09-02}{Bearing Sets the Price}{People treat you at the level you act --- until you pass the ego of whoever is above you}

\begin{core}
Treatment follows bearing. Ask and act at a level, and others supply the reasons for the confidence; the limits you fail to display are the limits you are not assigned. The same mechanism has a ceiling: a superior's self-image is an asset too, and visible brilliance is a claim against it. Outshine the master and the master survives you --- pleasantly, with a blander replacement already chosen.
\end{core}
\begin{mechanism}
Self-belief radiates like a crown. Observers cannot see the interior, so they infer from the exterior that there must be grounds: confidence is treated as evidence of competence because checking is expensive. Crown-bearers show no sense of limits, and limits disappear from what is asked and what is granted. Rebuff trains shrinking: the adult apologises for requests, pricing them small. Upward, the same display runs against insecurity. Superiors want to feel secure and superior; talent displayed beside them reads as a claim on their chair. The response is rarely open: appreciation is feigned, and at the first opportunity the sparkling one is replaced with someone less threatening. It fails two ways --- by boasting, a direct claim; and inadvertently, because some masters are so insecure that unguarded excellence alone offends.
\end{mechanism}
\begin{tells}
  \item Apologetic framing: pre-excuses, requests made small so rejection is cheap.
  \item A superior who praises warmly and blocks quietly.
  \item The bland promotion: the brilliant one passed over for someone non-threatening.
  \item Confidence with no visible basis, being rewarded --- the room is supplying the reasons.
  \item A colleague made visibly smaller in meetings their own work should have made large.
\end{tells}
\begin{moves}
  \item Set the asking level before the negotiation, not during it, and remove apology from requests entirely.
  \item Deliver quality quietly and attach the credit upward. Make the master look more brilliant than they are; the debt is repaid in protection.
  \item Never boast to a superior. The claim is the offence; competence delivered under their flag is not.
  \item Diagnose insecurity early: what happened to the last excellent person under them?
  \item Where the master cannot be soothed, route around them or leave; bearing cannot win an ego war it is not permitted to fight.
\end{moves}
\begin{counters}
  \item If you are the superior: price ideas, not deference. The comfortable candidate is not the safe one --- comfort was why they were chosen.
  \item Distinguish confidence from arrogance by behaviour under correction, not by volume.
  \item Watch what a hierarchy rewards. A hall of scraped bows is a hall of shrinking ideas.
  \item If you are the one shrinking, the limits are self-imposed: less was asked, so less was given.
  \item Test the grounds behind someone's bearing rather than the bearing itself (T-23-03), and separate the person from the trappings of office (T-18-01).
\end{counters}
\begin{cost}
The line between confidence and arrogance is drawn by the audience, not the actor, and misreading the season is expensive: regal stiffness in a reforming mood gathers revolt, and the thrones lost to it were lost with their heads. Loom too high above the crowd and you are the easiest target on the field. Elevating yourself by making others small forfeits the mechanism entirely --- height bought that way is priced as theft.
\end{cost}
\begin{field}
Any hierarchy: employment, institutions, courts, client work. Sharpest where one person controls your access to resources and cannot be cheaply replaced, because there the ceiling is a single ego.
\end{field}

\sources{B3}
\seesources
\seealso{T-04-03, T-18-01, T-23-03}
